\section{Exercise 1: Delay}
\subsection{Task and method}
Implement \texttt{[yn, yorigin] = delay(xn, xorigin, k)} for
\begin{equation}
  y(n)=x(n-k),\qquad k>0,
\end{equation}
where $k$ is an integer number of samples. A sample originally at $n_i$
appears at $n_i+k$ in the output. The vector values and their order are
unchanged. The output origin follows from
\begin{align}
  i-\texttt{yorigin} &= i-\texttt{xorigin}+k,\\
  \texttt{yorigin} &= \texttt{xorigin}-k.
\end{align}
The implementation constructs the two time axes separately and displays
the signals in two subplots within one figure.

\subsection{Source code}
\lstinputlisting[caption={Scilab delay function and example call},label={lst:delay}]{../1.sci}

\subsection{Example and expected result}
For the example in the lab handout,
\begin{lstlisting}[numbers=none]
[yn, yorigin] = delay([1, -2, 3, 6], 3, 1);
\end{lstlisting}
the calculated output is
\[
  \texttt{yn}=[1,-2,3,6],\qquad \texttt{yorigin}=2.
\]
\begin{table}[H]
  \centering
  \caption{Sample coordinates before and after a one-sample delay}
  \label{tab:delay}
  \begin{tabular}{lrrrr}
    \toprule
    Vector position $i$ & 1 & 2 & 3 & 4\\
    \midrule
    Sample value & 1 & $-2$ & 3 & 6\\
    Input time $n_i$ & $-2$ & $-1$ & 0 & 1\\
    Output time $n_i+1$ & $-1$ & 0 & 1 & 2\\
    \bottomrule
  \end{tabular}
\end{table}

\subsection{Execution evidence}
\resultfigure{images/exercise1.png}
  {Input and delayed output for Exercise 1}
  {fig:delay}
The captured Scilab plot agrees with the calculated sample coordinates
in \Cref{tab:delay}; the output subplot contains only the delayed signal.

\subsection{Discussion}
The captured output graph is shifted to the right by one sample. In particular,
the sample of amplitude $3$ moves from $n=0$ to $n=1$, while the sample
of amplitude $6$ moves from $n=1$ to $n=2$. The output value at $n=0$ is
$-2$, which occupies vector position two. This explains why the origin
parameter decreases even though the signal moves to the right.
